% =====================================================================
%  CARNET D'ENTRAINEMENT — FICHE 1 : Les chiffres significatifs
%  THÈME 3 — Arrondir & opérations (addition/soustraction, mult./division)
%  Fiche TUTO
% =====================================================================
\documentclass[11pt,a4paper]{article}
\usepackage[utf8]{inputenc}
\usepackage[T1]{fontenc}
\usepackage{lmodern}
\usepackage[french]{babel}
\usepackage{amsmath,amssymb,mathtools}
\usepackage{icomma}
\usepackage{siunitx}
\sisetup{output-decimal-marker={,},per-mode=symbol}
\usepackage[version=4]{mhchem}
\usepackage{xcolor}
\usepackage{tikz}
\usepackage[most]{tcolorbox}
\usepackage{enumitem}
\usepackage{booktabs}
\usepackage{array}
\usepackage{fancyhdr}
\usepackage[a4paper,margin=2cm,headheight=26pt,headsep=8pt]{geometry}
\usetikzlibrary{arrows.meta,calc}

\definecolor{primary}{RGB}{150,98,162}
\definecolor{primarydark}{RGB}{92,52,110}
\definecolor{reglecol}{RGB}{150,98,162}
\definecolor{methcol}{RGB}{90,90,100}
\definecolor{excol}{RGB}{0,135,90}
\definecolor{attncol}{RGB}{200,40,55}
\newcommand{\cs}{chiffre significatif}
\newcommand{\CS}{chiffres significatifs}

\pagestyle{fancy}\fancyhf{}
\fancyhead[L]{\small\textcolor{primarydark}{\textbf{Fiche 1} — Chiffres significatifs}}
\fancyhead[R]{\small\textcolor{primarydark}{\textsc{Tuto · Thème 3}}}
\fancyfoot[C]{\small\thepage}
\renewcommand{\headrulewidth}{0.8pt}
\renewcommand{\headrule}{\hbox to\headwidth{\color{primary}\leaders\hrule height \headrulewidth\hfill}}

\tcbset{boxbase/.style={enhanced,breakable,boxrule=0.9pt,arc=2.5pt,
  left=3mm,right=3mm,top=1.6mm,bottom=1.6mm,fonttitle=\bfseries,coltitle=white}}
\newtcolorbox{regle}[1][]{boxbase,colback=reglecol!7,colframe=reglecol,
  title={\textsc{Règle}\ifx\relax#1\relax\relax\else\ ---\ \textmd{\itshape #1}\fi}}
\newtcolorbox{methode}[1][]{boxbase,colback=methcol!7,colframe=methcol,sharp corners,
  title={\textsc{Méthode}\ifx\relax#1\relax\relax\else\ : \textmd{\itshape #1}\fi}}
\newtcolorbox{exemple}[1][]{boxbase,colback=excol!6,colframe=excol,
  title={\textsc{Exemple}\ifx\relax#1\relax\relax\else\ : \textmd{\itshape #1}\fi}}
\newtcolorbox{attention}[1][]{boxbase,colback=attncol!6,colframe=attncol,
  title={\textsc{Attention}\ifx\relax#1\relax\relax\else\ : \textmd{\itshape #1}\fi}}

\setlength{\parindent}{0pt}\setlength{\parskip}{2pt}

\begin{document}

\begin{center}
{\LARGE\bfseries\color{primarydark} Thème 3 — Arrondir \& opérations}\\[1pt]
{\large\color{primary} Le résultat n'est jamais plus précis que les données}
\end{center}
\vspace{1mm}

Un calcul ne peut pas \emph{fabriquer} de la précision : le résultat doit refléter la donnée la moins
bien connue. Ce thème entraîne trois gestes : \textbf{arrondir} proprement, puis appliquer la
\textbf{bonne règle} selon qu'on additionne ou qu'on multiplie.

\begin{regle}[arrondir au rang voulu]
On regarde le \textbf{premier chiffre supprimé} (juste à droite du dernier chiffre gardé) :
$0,1,2,3,4 \rightarrow$ on arrondit \textbf{par défaut} (chiffre gardé inchangé) ;
$5,6,7,8,9 \rightarrow$ on arrondit \textbf{par excès} ($+1$ au chiffre gardé).
Une \emph{retenue} peut se propager : $\num{2,996}$ à $3$~CS donne $\num{3,00}$ (on garde le
« $00$ » pour conserver $3$~CS).
\end{regle}

\begin{regle}[addition et soustraction : on compte les \emph{décimales}]
Le résultat s'arrondit au \textbf{même nombre de décimales} (chiffres après la virgule) que la donnée
qui en a le \emph{moins}. (Une somme se fait rang par rang : le rang le moins bien connu contamine ce
même rang du résultat.)
\end{regle}

\begin{regle}[multiplication et division : on compte les \CS]
Le résultat s'arrondit au \textbf{même nombre de \CS} que la donnée qui en a le \emph{moins}.
\end{regle}

\begin{exemple}[chaque règle sur un cas]
\textbf{Addition} — la donnée la moins précise a $1$ décimale :
\[ \num{123,25}+\num{46,0}+\num{86,257}=\num{255,507}\ \xrightarrow{\text{1 décimale}}\ \boxed{\num{255,5}}. \]
\textbf{Division} — la donnée la moins riche a $3$~CS :
\[ \frac{\qty{0,0361}{\milli\gram}}{\qty{12,0}{\litre}} = \num{3,00833e-3}\ \xrightarrow{3\ \text{CS}}\ \boxed{\num{3,01e-3}}\ \si{\gram\per\cubic\metre}. \]
\end{exemple}

\begin{attention}[deux réflexes de survie]
\textbf{1. On n'arrondit qu'à la fin.} Dans un calcul en plusieurs étapes, on garde des
\emph{chiffres de garde} et on arrondit \textbf{une seule fois}, au résultat final. Arrondir trop tôt
amplifie l'erreur.

\smallskip
\textbf{2. La soustraction de nombres voisins détruit les \CS.} $\num{8,11}-\num{8,0}=\num{0,11}\to\num{0,1}$ :
de $3$ et $2$~CS, on tombe à $\mathbf{1}$~CS. À éviter en chimie analytique.
\end{attention}

\begin{methode}[un calcul mixte $+$ et $\times$]
\begin{enumerate}[leftmargin=1.6em,topsep=1pt,itemsep=0pt]
  \item faire tout le calcul \emph{sans arrondir} (chiffres de garde) ;
  \item repérer, pour \emph{chaque} étape, la règle adaptée (décimales pour $+/-$, CS pour $\times/\div$) ;
  \item arrondir \emph{une seule fois}, à la fin, selon l'étape limitante.
\end{enumerate}
\end{methode}

\begin{center}\small
\renewcommand{\arraystretch}{1.15}
\begin{tabular}{@{}l l l@{}}
\toprule
\textbf{Opération} & \textbf{On compte\dots} & \textbf{Le résultat a\dots}\\
\midrule
Addition / soustraction   & les décimales & le minimum de décimales des données\\
Multiplication / division & les \CS       & le minimum de \CS\ des données\\
\bottomrule
\end{tabular}
\end{center}

\end{document}
